\section{Limitations and Conclusion}
\label{sec:limitations}
\label{sec:limitations-conclusion}

The shared-threshold law is a boundary-motion model, not a posterior identified
by calibrated marginals. More quadrature nodes reduce numerical error, not
posterior discrepancy. The nHD theorem states a sufficient surface-geometric
condition; one label per image cannot verify it, and observed support failures
show where it can break.

Full nHD is sensitive to one extreme surface point. Pooled nHD95 is more robust
but nonmetric and empirically sensitive to checkpoint and action threshold.
Surface extraction, spacing, normalization, pooling, and empty-mask penalties
must therefore be fixed as deployment choices. Our five binary datasets, two
model families, primary seed, and image bootstrap do not cover site,
multi-rater, patient-cluster, or training uncertainty. Learned detectors,
calibrated annotation posteriors, multiclass actions, and absolute risk
calibration remain future work; Appendix~\ref{app:limitations} gives details.

\paragraph{Conclusion.}
\label{sec:conclusion}
Selective segmentation begins by naming the failure that should trigger
abstention. Loss indexing supplies the target, but the working posterior must
also encode how that failure occurs. Shared thresholds model coherent boundary
motion, and normalized full Hausdorff loss yields the clearest theory and
empirical ranking. nHD95 exposes the value and sensitivity of a robust endpoint;
Dice explains why the same coupling need not dominate count-oriented SDC.
Confidence is not universal or metric-free: it is a declared model of how a
fixed deployed mask may be wrong.
